%%============================================================================
%% A report. Created by report-new.sh.
%% Halil asked on 2026-11-13 for this report without the canonical sections, so
%% it is compiled by hand in a temporary folder, the way report-build.sh does,
%% and report-check.py is not run on it.
%%============================================================================
\documentclass{labreport}
\lrnoreflection   % lib/reports/src/config.yml says reflection: false

\lrdate{2026-11-13}
\lrslug{halil-ideal-functionalities-and-attacks}
\lrtitle{Ideal functionalities and the attacks they admit}
\lrauthor{Halil İbrahim Kanpak}
\lrstatus{draft}

%% Every number in prose is defined here, with the file it came from.
\lrdefine{AttackModels}{103}{lib/reports/src/data/2026-11-13-halil-ideal-functionalities-and-attacks-facts.csv}
\lrdefine{NeedUpdate}{40}{lib/reports/src/data/2026-11-13-halil-ideal-functionalities-and-attacks-facts.csv}
\lrdefine{NeedRound}{30}{lib/reports/src/data/2026-11-13-halil-ideal-functionalities-and-attacks-facts.csv}
\lrdefine{NeedMore}{70}{lib/reports/src/data/2026-11-13-halil-ideal-functionalities-and-attacks-facts.csv}
\lrdefine{NeedCut}{12}{lib/reports/src/data/2026-11-13-halil-ideal-functionalities-and-attacks-facts.csv}
\lrdefine{NeedFinal}{4}{lib/reports/src/data/2026-11-13-halil-ideal-functionalities-and-attacks-facts.csv}
\lrdefine{NeedDecrypted}{4}{lib/reports/src/data/2026-11-13-halil-ideal-functionalities-and-attacks-facts.csv}
\lrdefine{NeedRest}{13}{lib/reports/src/data/2026-11-13-halil-ideal-functionalities-and-attacks-facts.csv}
\lrdefine{SalemGames}{18}{lib/reports/src/data/2026-11-13-halil-ideal-functionalities-and-attacks-facts.csv}

\newcommand{\IF}{\mathcal{F}}
\newcommand{\IG}{\mathcal{G}}

\begin{document}
\lrmaketitle

%%----------------------------------------------------------------------------
\section{What we have}

A setting of learning is an \emph{ideal functionality}, a \emph{scenario} is an ideal functionality with a coalition and the coalition's behavior, and an \emph{attack} is a question asked in a scenario.
The possible attacks come from the two sides of an ideal functionality that face the coalition, what it hands out and what it accepts, so they are generated from the settings and not collected from papers.
Ten standard attacks result, in three question kinds, and the library holds ten ideal functionalities, among them FedAvg in three variants, split and vertical learning, private inference and HE-OFT.
Two ideal functionalities with the same parties are ordered by what they hand a coalition, and an attack's advantage can only grow along that order (proven).
Each attack is therefore located by its lowest ideal functionalities, so the relation between attacks and ideal functionalities is a family of up-sets and not an arbitrary many-to-many map.
The lowest ideal functionalities of an attack say where encryption must go to stop it, and the floor of a task, the lowest ideal functionality that still computes it, holds the attacks no encryption stops.
How much an attack gains where it runs is a separate quantitative question, answered by a bound or by a measurement.

\lrboundary{The definitions are ours, written in the simulation paradigm that secure aggregation is already proven in, and no other reader has checked them. An empirical verdict rests on the ledger's coding of the paper, and of the works cited only Yeom et al., Salem et al. and HE-OFT were read in full.}

%%----------------------------------------------------------------------------
\section{Scenarios}

The parties are a server, clients $1,\dots,n$ and a querier, and each party $P$ holds an input $x_P$.
Client $i$ holds a dataset $D_i$, a multiset of records, the server a model or the messages it chooses, and the querier a query.
For a set $B$ of parties, $x_B$ collects their inputs.

An ideal functionality $\IF$ is one incorruptible party that receives every input and hands each party $P$ an output $\IF_P(x)$.
It also fixes an \emph{input space} for each party, the set of inputs it accepts from that party.
A setting that runs in rounds takes a fresh input each round and hands out a sequence of outputs, and $\IF_P(x)$ then denotes the whole sequence.
The outputs are the side the coalition observes and the input spaces the side through which it acts, and the library below writes both for every setting.

A \emph{coalition} $A$ is the set of parties the adversary controls, and $H$ is the set of the honest parties.
The coalition's \emph{view} under $\IF$ is $(x_A,\IF_A(x))$, its own inputs and every output handed to one of its members.
Its \emph{behavior} is one of three.
\begin{itemize}
  \item A \emph{semi-honest} coalition supplies its true inputs and only observes.
  \item An \emph{input-choosing} coalition supplies any inputs its input spaces allow, chosen from what it has seen so far, and aborts where $\IF$ permits it. A protocol secure against malicious coalitions leaves the adversary exactly this power.
  \item A \emph{deviating} coalition departs from the protocol. Against a protocol secure only for semi-honest coalitions the simulation then says nothing, so such an attack is placed on the protocol and not on $\IF$.
\end{itemize}
A scenario is the triple $(\IF,A,b)$ of an ideal functionality, a coalition and a behavior $b$.

A protocol $\pi$ \emph{realizes} $\IF$ against $A$ with error $\epsilon$ when a simulator, given only $(x_A,\IF_A(x))$, produces the coalition's view in $\pi$ up to distance $\epsilon$.
The distance is taken jointly with the honest outputs and holds for every input $x$, which is the definition secure aggregation is proven against~\cite{MOJC23,MWAPR23}.
Every question below then gains at most $2\epsilon$ more in $\pi$ than under $\IF$ (proven, since each hypothesis moves the coalition's view by at most $\epsilon$).
A question therefore splits into a cryptographic layer, the size of $\epsilon$, and an output layer, the advantage under $\IF$, and the rest of this report concerns the output layer.
Threshold CKKS computes the outputs only approximately, so it realizes a variant of $\IF$ whose outputs carry the CKKS error, and its realization error rests on a smudging hypothesis (heuristic, \lrpath{lib/reports/2026-11-11-halil-he-results-carried-over.pdf}).

%%----------------------------------------------------------------------------
\section{Questions and the standard attacks}

A question is posed on the honest inputs $x_H$ and the honest outputs, and it comes in one of three kinds.
\begin{itemize}
  \item In a \emph{distinguish} question two distributions $\mathcal{D}_0$ and $\mathcal{D}_1$ over $x_H$ are fixed, $x_H$ is drawn from $\mathcal{D}_\beta$ for a hidden bit $\beta$, and the coalition guesses $\beta$ from its view. Its advantage is $\lvert\Pr[\text{guess } 1\mid\beta=1]-\Pr[\text{guess } 1\mid\beta=0]\rvert$.
  \item In an \emph{invert} question a target function $g$, a loss $\ell$ and a tolerance $\tau$ are fixed, and the coalition wins when its output $y$ satisfies $\ell(y,g(x_H))\le\tau$. Its advantage is its probability of winning minus that of the best guess made from the prior on $x_H$ alone.
  \item In a \emph{steer} question a target set $T$ of honest outputs is fixed, and an input-choosing coalition picks inputs $x'_A$ from its input spaces. Its advantage is $\Pr[\IF_H(x'_A,x_H)\in T]$ minus the same probability under its true inputs.
\end{itemize}
Distinguish and invert questions act through what $\IF$ hands out and steer questions through what it accepts, so a semi-honest coalition has zero steering advantage by definition.

A standard attack is one question kind with one target that the input and output types of the library admit, and Table~\ref{attacks} lists the ten.
The names are targets, so membership inference in a paper is the membership question asked in that paper's scenario.
Source keeps the union and every dataset size fixed, because an ideal functionality that hands out the sizes would otherwise answer it outright.
By construction every honest input type has an invert target and every honest output a steer target.
Whether the list covers every attack of interest is open, and free-riding is a known case outside it, since it concerns what the coalition contributes.

% Generated by mktable.py from lib/reports/src/data/2026-11-13-halil-ideal-functionalities-and-attacks-attacks.csv
% Do not edit. Regenerate: see lab/analysis/
% Requires: \usepackage{booktabs}
\begin{table}[!htbp]
  \centering
  \caption{The standard attacks. Each is one question kind with one target, and the interface is the side of the ideal functionality it acts through.}
  \label{attacks}
  \small\begin{tabular}{>{\raggedright\arraybackslash}p{0.13\linewidth} >{\raggedright\arraybackslash}p{0.12\linewidth} >{\raggedright\arraybackslash}p{0.53\linewidth} >{\raggedright\arraybackslash}p{0.10\linewidth}}
    \toprule
    \textbf{Name} & \textbf{Kind} & \textbf{The two hypotheses, or the target} & \textbf{Interface} \\
    \midrule
    membership & distinguish & a record $z$ is in an honest dataset, or it is not & outputs \\
    source & distinguish & $z$ is with honest client $i$ and $z'$ with honest client $j$, or the two are swapped, so the union and every size stay fixed & outputs \\
    property & distinguish & the honest data come from a distribution with a property, or without it & outputs \\
    record & invert & a record of an honest dataset, within a distance & outputs \\
    attribute & invert & one feature or the label of a record whose other features are known & outputs \\
    model & invert & the parameters or the behavior of a model that is an honest input & outputs \\
    query & invert & the query of an honest querier & outputs \\
    targeted & steer & honest outputs whose model maps chosen inputs to chosen labels & inputs \\
    untargeted & steer & honest outputs whose model has accuracy below a level & inputs \\
    delivery & steer & the honest parties receive no output & inputs \\
    \bottomrule
  \end{tabular}
\end{table}


%%----------------------------------------------------------------------------
\section{The library of ideal functionalities}

Table~\ref{library} writes each setting as an ideal functionality, with what it accepts and what it hands to the server, to a client and to the querier.
Variants of one setting are different ideal functionalities, and fedavg-plain, fedavg-secagg and fedavg-final differ only in what the server receives.
The server's input in fedavg-plain and fedavg-secagg is the model it sends each client each round, because deployed protocols let the server choose it.
Requiring one model for all clients gives a variant with a smaller input space, which a protocol realizes only by checking that every client received the same model.
The clients' local training is part of $\IF$, so an honest client gives its dataset and $\IF$ computes its update, while an input-choosing client may submit any update.
In he-oft the head displacement and the per-class counts are computed in the same way, and the server's output is the leakage of the simulation theorem of HE-OFT~\cite{KSK26}.
The ideal functionality outsourced-training has a single data owner, so the framework also covers settings that are not multi-party.

% Generated by mktable.py from lib/reports/src/data/2026-11-13-halil-ideal-functionalities-and-attacks-library.csv
% Do not edit. Regenerate: see lab/analysis/
% Requires: \usepackage{booktabs}
\begin{table}[!htbp]
  \centering
  \caption{The library of ideal functionalities. A client that chooses its input may submit anything the Accepts column allows, and every output not listed is withheld.}
  \label{library}
  \small\begin{tabular}{>{\raggedright\arraybackslash}p{0.14\linewidth} >{\raggedright\arraybackslash}p{0.26\linewidth} >{\raggedright\arraybackslash}p{0.19\linewidth} >{\raggedright\arraybackslash}p{0.15\linewidth} >{\raggedright\arraybackslash}p{0.14\linewidth}}
    \toprule
    \textbf{Name} & \textbf{Accepts} & \textbf{The server receives} & \textbf{A client receives} & \textbf{The querier receives} \\
    \midrule
    pooled-training & a dataset from each client & the model trained on the shuffled union & the same model & nothing \\
    served-model & a dataset from each client, queries & nothing & nothing & the trained model's answer to each query \\
    outsourced-training & a dataset from one client & nothing & the model trained on it & nothing \\
    fedavg-plain & a dataset or any update from each client, a model for each client each round from the server & every client's update in every round & the model sent to it & nothing \\
    fedavg-secagg & as fedavg-plain & the sum of the updates in every round & the model sent to it & nothing \\
    fedavg-final & a dataset or any update from each client & the final model & the average model of every round & nothing \\
    split-learning & batches from the client, any gradient at the cut from the server & the activations at the cut and the labels & the gradient at the cut & nothing \\
    vertical-training & a feature block of the shared records from each client, the labels and any gradient from the server & every client's embedding & the gradient at its own embedding & nothing \\
    private-inference & a model from the server, a query & nothing & nothing & the model's answer to the query \\
    he-oft & a dataset, or any head displacement and per-class counts, from each client, up to $Q$ queries~\cite{KSK26} & each client's sample count and the selected arrangement & the selected arrangement & one label per query \\
    \bottomrule
  \end{tabular}
\end{table}


%%----------------------------------------------------------------------------
\section{The order, and where an attack runs}

Fix a coalition $A$ and two ideal functionalities $\IF$ and $\IG$ with the same parties and the same honest input spaces.
$\IF$ is \emph{below} $\IG$ for $A$, written $\IF\preceq_A\IG$, when an efficient map $m$ satisfies
\[
  \bigl(x_A,\IF_A(x)\bigr)\;\stackrel{d}{=}\;m\bigl(x_A,\IG_A(x)\bigr)\quad\text{for every input } x,
\]
so that the coalition computes its view under $\IF$ from its view under $\IG$.
If $\IF\preceq_A\IG$, a semi-honest coalition gains no more under $\IF$ than under $\IG$ in any distinguish or invert question (proven, since a strategy under $\IF$ runs under $\IG$ after $m$).
For a level $\delta$, the ideal functionalities on which an attack's advantage exceeds $\delta$ thus form an up-set, and the attack is located by the minimal elements of that set.
Steering is ordered by the input spaces instead, because an ideal functionality that computes the same honest outputs and accepts fewer inputs from $A$ admits no more steering (proven, as a maximum over a smaller set).

For a semi-honest server, fedavg-final $\preceq$ fedavg-secagg $\preceq$ fedavg-plain (proven), because the server computes the final model from the round sums and each round sum from the updates.
The ideal functionality pooled-training is comparable with none of the three, because its output does not depend on how the records are split among the clients and theirs does.

An input-choosing server acts through both sides at once and can move itself up the order.
Pasquini et al.~\cite{PFA22} send a crafted model to every client but the target, so the round sum under fedavg-secagg equals the target's update, which is what fedavg-plain hands out (empirical).
The variant with one model for all clients blocks that route, yet Shi et al.~\cite{SWXZSHL25} still reconstruct the records of the union with one crafted model, without attributing them to clients (empirical).
In the terms of Table~\ref{attacks}, the variant still admits record, and that attack gives no answer to source.
The order therefore locates an active attack only once the input space is fixed, so the input space is a coordinate of an ideal functionality as much as its outputs are.

%%----------------------------------------------------------------------------
\section{Where encryption goes, and the floor}

A protocol stops an attack in a scenario when the ideal functionality it realizes lies outside the attack's up-set.
Encryption and MPC let a protocol realize a lower ideal functionality of the same task, fedavg-secagg in place of fedavg-plain for instance, and each step down costs computation.
Placing HE is then the choice of the highest ideal functionality of the task, the one with the least encryption, that lies outside the up-set of every attack the deployment must stop.

A task fixes what some party must receive, such as the trained model for the clients or one label per query for the querier.
Its \emph{floor} is the lowest ideal functionality of the library that still hands that out, and an attack whose up-set contains the floor runs under every protocol for the task.
Encryption cannot stop such an attack, and what remains is to change the ideal functionality, by DP inside it or by narrowing what it accepts, or to record the leakage as inherent.
Each attack in a scenario thus gets one of three verdicts at the output layer: zero advantage, where encryption alone suffices, a bound from DP inside the ideal functionality, or no bound, where a measurement or a gap is recorded.
Table~\ref{cases} gives worked cases.

Source has zero advantage under fedavg-secagg with one full-batch local step weighted by size, because client $i$'s update at round $t$ is the sum of the per-record gradients of $D_i$, so the round sum is
\[
  \sum_{i=1}^{n}\sum_{z\in D_i}\nabla L(\theta_t,z)\;=\;\sum_{z\in D_1\uplus\cdots\uplus D_n}\nabla L(\theta_t,z),
\]
a function of the union (proven), where $\theta_t$ is the model of round $t$, $L$ the training loss and $\uplus$ the union of multisets.
Several local steps break the identity, because the update is then a nonlinear function of the client's own records, so the round sum depends on the split.

HE-OFT shows the floor of a task that serves labels.
Its shared displacement for class $c$ is the count-weighted mean of the clients' displacements $\Delta_{j,c}$ with per-class counts $g_{j,c}$, so a client $a$ that inflates its count moves it at will,
\[
  \frac{\sum_{j} g_{j,c}\,\Delta_{j,c}}{\sum_{j} g_{j,c}}\;\longrightarrow\;\Delta_{a,c}\qquad\text{as } g_{a,c}\to\infty
\]
(proven, prop-he-oft-count-steering of the carried-over report).
The counts are inputs he-oft accepts, so encryption leaves this steering in place, and only a variant that bounds the accepted counts removes it.

% Generated by mktable.py from lib/reports/src/data/2026-11-13-halil-ideal-functionalities-and-attacks-cases.csv
% Do not edit. Regenerate: see lab/analysis/
% Requires: \usepackage{booktabs}
\begin{table}[!htbp]
  \centering
  \caption{Worked cases. The verdict is about the ideal functionality, so a protocol realizing it adds at most its realization error. The active attacks on fedavg-secagg are discussed in the text.}
  \label{cases}
  \small\begin{tabular}{>{\raggedright\arraybackslash}p{0.12\linewidth} >{\raggedright\arraybackslash}p{0.27\linewidth} >{\raggedright\arraybackslash}p{0.40\linewidth} >{\raggedright\arraybackslash}p{0.09\linewidth}}
    \toprule
    \textbf{Attack} & \textbf{Scenario} & \textbf{Verdict at the output layer} & \textbf{Status} \\
    \midrule
    source & pooled-training, any coalition holding neither $z$ nor $z'$ & zero, because every output is a function of the union and the coalition's inputs & proven \\
    source & fedavg-secagg or fedavg-final, semi-honest server, one full-batch local step weighted by size & zero, because every round's sum is the sum of the per-record gradients over the union & proven \\
    source & the same, with several local steps & the model depends on how the records are split, so the advantage need not vanish, and its size is unknown & open \\
    source & fedavg-plain, semi-honest server & at least the advantage of telling whether one client's update was computed with $z$ or with $z'$, since the server reads that update alone & proven \\
    membership & every training task that hands out a model or answers & in the floor, and at most $e^\varepsilon-1$ when the ideal functionality trains with $\varepsilon$-DP~\cite{YGFJ18} & proven \\
    targeted, untargeted & pooled-training, input-choosing client & in the floor, because the ideal functionality accepts any dataset & proven \\
    targeted & he-oft, input-choosing client & in its floor, because inflated per-class counts move a row of the shared head at will~\cite{KSK26} & proven \\
    membership & he-oft, querier with $Q$ queries & in its floor, and at most the advantage of a holder of both heads in plaintext~\cite{KSK26} & proven \\
    every output question & outsourced-training, server coalition & zero, because the server receives nothing & proven \\
    \bottomrule
  \end{tabular}
\end{table}


%%----------------------------------------------------------------------------
\section{Published attacks, placed}

The ledger codes, for each of its \lrv{AttackModels} attack threat models, what the attack reads in plaintext, and the least of it gives the lowest ideal functionality of the library under which the attack runs as published (empirical, coded on 2026-10-16 and not re-read).
Of these, \lrv{NeedUpdate} read a client's own update and \lrv{NeedRound} a per-round sum or model, so an ideal functionality that hands the server only the final model stops these \lrv{NeedMore} as published.
Another \lrv{NeedCut} read values at a cut, which places them under split-learning and vertical-training.
Only \lrv{NeedFinal} need no more than the final model or served answers, which puts them in the floor of every training task, and \lrv{NeedDecrypted} read a decrypted value, which attacks the realization and not the ideal functionality.
The other \lrv{NeedRest} are coded as reading nothing in plaintext or values outside the library, and the placement ignores the input space an attack uses.

%%----------------------------------------------------------------------------
\section{Prior work}

Table~\ref{prior} lists the works this report compares against, what each fixes in the picture above, and where its PDF is held.
Each fixes one piece, the vocabulary, the games on one model, the simulation of one protocol or one joint construction, and the gap they leave is the subject of \lrpath{lib/reports/2026-11-11-halil-prior-work-on-scenario-frameworks.pdf}.
Barreno et al. and the three standards are not held, the first behind the Springer paywall and the others because they are sold.

% Generated by mktable.py from lib/reports/src/data/2026-11-13-halil-ideal-functionalities-and-attacks-prior.csv
% Do not edit. Regenerate: see lab/analysis/
% Requires: \usepackage{booktabs}
\begin{table}[!htbp]
  \centering
  \caption{The prior work this report compares against, what each fixes in the picture of ideal functionalities and attacks, and the folder of its PDF under lib/literature/pdfs, where the file name starts with the key.}
  \label{prior}
  \small\begin{tabular}{>{\raggedright\arraybackslash}p{0.24\linewidth} >{\raggedright\arraybackslash}p{0.55\linewidth} >{\raggedright\arraybackslash}p{0.12\linewidth}}
    \toprule
    \textbf{Work} & \textbf{What it fixes in this picture} & \textbf{PDF folder} \\
    \midrule
    Barreno et al.~\cite{BNJT10} & goal and capability, which become the two interfaces & not held \\
    Papernot et al.~\cite{PMSW18} & the stage of an attack and the adversary's knowledge & ieee-eurosp \\
    NIST AI 100-2~\cite{VOFADH25} & the names of attacks, privacy and integrity together & nist \\
    Yeom et al.~\cite{YGFJ18} & membership and attribute games on one model, the DP bound & csf \\
    Salem et al.~\cite{SCEKPSTZ23} & \lrv{SalemGames} games on one model, the output layer of pooled-training & ieee-sp \\
    Li et al.~\cite{LLLKPMW24} & games on gradients, the output layer of fedavg-plain & acm-ccs \\
    Wu et al.~\cite{WZLLHL23} & integrity attacks, the steer kind, with no advantage & arxiv \\
    Mansouri et al.~\cite{MOJC23} & the ideal functionality of secure aggregation, by coalition & popets \\
    Ma et al.~\cite{MWAPR23} & simulation of secure aggregation over many rounds & ieee-sp \\
    Morona-M\'{\i}nguez et al.~\cite{MPP26} & averaging under multi-key HE, up to the CKKS error & arxiv \\
    HE-OFT~\cite{KSK26} & an ideal functionality with stated leakage, and the bound on label queries & arxiv \\
    Pentyala et al.~\cite{PRMMDND22} & MPC for the protocol and DP inside the ideal functionality & arxiv \\
    Carletti et al.~\cite{CFMPV25} & the threat models of gradient inversion, by role & usenix-security \\
    Rizzini et al.~\cite{REGB26} & verifiable integrity, which narrows what is accepted & popets \\
    Desfontaines and Pej\'o~\cite{DP20} & criteria for a list of attacks, exclusive and compatible dimensions & popets \\
    NIST SP 800-226~\cite{NISTSP800226} & the unit of privacy for DP inside an ideal functionality & nist \\
    Villa et al.~\cite{VDKMR26} & one game for confidentiality, integrity and availability of agentic LLM systems & arxiv \\
    IEEE 3652.1~\cite{IEEE3652}, ISO/IEC 4922~\cite{ISO4922}, ISO/IEC 28033~\cite{ISO28033} & the vocabulary of federated learning, MPC and FHE & not held \\
    \bottomrule
  \end{tabular}
\end{table}


%%----------------------------------------------------------------------------
\section{Open}

\begin{itemize}
  \item Active coalitions need an order on pairs of outputs and input spaces, and the up-set statement is proven here only for a fixed input space.
  \item Settings that run in rounds against adaptive coalitions need the order round by round, which the definition on whole output sequences does not give.
  \item Each ideal functionality should be stated with the noisy output that threshold CKKS computes, so that the realization error concerns the protocol alone.
  \item The library lacks decentralized learning, fine-tuning with LoRA adapters across parties, and inference with a model split across parties.
  \item Source under several local steps is open, and it decides whether the example of client A and client B holds for FedAvg as deployed.
\end{itemize}


\lrreferences

\end{document}
